\documentclass[a4paper,10.5pt]{article}
\usepackage[utf8]{inputenc}
\usepackage[T1]{fontenc}
\usepackage[french]{babel}
\usepackage[left=1cm,right=1cm,top=.5cm,bottom=0cm]{geometry}
\usepackage{enumitem}
\usepackage{hyperref}
\usepackage{xcolor}
\usepackage{titlesec}
\usepackage{fontawesome5}
\usepackage{lmodern}
\usepackage[normalem]{ulem}

% --- Couleurs Safran ---
\definecolor{primary}{RGB}{0, 51, 102}
\definecolor{secondary}{RGB}{80, 80, 80}

% --- Configuration des liens ---
\hypersetup{colorlinks=true, linkcolor=primary, filecolor=primary, urlcolor=primary}

% --- Formatage des titres ---
\titleformat{\section}{\large\bfseries\color{primary}\uppercase}{}{0em}{}[\titlerule]
\titlespacing{\section}{0pt}{8pt}{5pt}

% --- Commandes personnalisées ---
\newcommand{\resumeItem}[1]{\item\small{#1}}
\newcommand{\resumeSubheading}[4]{
  \vspace{-1pt}\item
    \begin{tabular*}{0.97\textwidth}[t]{l@{\extracolsep{\fill}}r}
      \textbf{#1} & \textit{\small #2} \\
      \textit{\small #3} & \textit{\small #4} \\
    \end{tabular*}\vspace{-5pt}
}

\begin{document}

% --- EN-TÊTE ---
\begin{center}
    {\Large \textbf{Loïc Aron MBASSI EWOLO}} \\ \vspace{4pt}
    \textbf{\large Stage : Déploiement Multi-Plateforme Modèles IA (GPU/NPU/FPGA)} \\
    \textit{\small Disponible dès \textbf{Avril 2026} pour 5 mois}\\
    \vspace{4pt}
    \small
    \faMapMarker\ Cergy, France (Mobile Massy / IDF) \quad
    \faPhone\ (+33) 7 59 52 33 98 \quad
    \faEnvelope\ \href{mailto:loicaron.mbassiewolo@ensea.fr}{loicaron.mbassiewolo@ensea.fr} \\ \vspace{2pt}
    \faLinkedin\ \href{https://www.linkedin.com/in/loic-mbassi}{loic-mbassi} \quad
    \faGithub\ \href{https://github.com/Nameless0l}{Nameless0l} \quad
    \faGlobe\ \href{https://mbassiloic.tech}{mbassiloic.tech}
\end{center}

\vspace{-6pt}

% --- PROFIL ---
\noindent\small{Élève-ingénieur en double diplôme \textbf{ENSEA/Polytechnique Yaoundé}, spécialisé en \textbf{systèmes embarqués} et \textbf{IA}. Expérience en déploiement de modèles \textbf{PyTorch} sur \textbf{GPU Nvidia} (Jetson), développement \textbf{FPGA} (VHDL/Quartus) et \textbf{TinyML} sur microcontrôleurs. Je recherche un stage chez \textbf{Safran} pour contribuer au déploiement multi-plateforme de modèles de vision embarquée.}

% --- FORMATION ---
\section{Formation}
\begin{itemize}[leftmargin=*, label={.}, nosep]
    \item \textbf{ENSEA (École Nationale Supérieure de l'Électronique et de ses Applications)} \hfill \textit{Cergy, France | 2025 -- Présent} \\
    \small{Ingénieur 2A, Double Diplôme - \textbf{Systèmes Embarqués}, Traitement du Signal et \textbf{IA}.} \\
    \textit{\small{Cours : Deep Learning, Computer Vision, Conception FPGA (VHDL), Architecture processeurs.}}
    \vspace{2pt}
    \item \textbf{École Nationale Supérieure Polytechnique} \hfill \textit{Yaoundé, Cameroun | 2021 -- 2025} \\
    \small{Diplôme d'Ingénieur de Conception (Niveau Master 2) : \textbf{Mathématiques Appliquées}, Génie Logiciel \& Systèmes.}
    \item \textbf{Université de Yaoundé I} \hfill \textit{Yaoundé, Cameroun | 2020 -- 2022} \\
    \small{Licence en Informatique et Mathématiques : Algorithmique C, Analyse, Algèbre Linéaire.}
\end{itemize}

% --- COMPÉTENCES TECHNIQUES ---
\section{Compétences Techniques}
\begin{itemize}[leftmargin=*, nosep]
    \item \small{\textbf{IA / Deep Learning :} \textbf{PyTorch}, TensorFlow, ONNX, Computer Vision (YOLO, OpenCV), TinyML.}
    \item \small{\textbf{Déploiement Embarqué :} \textbf{Nvidia Jetson} (TensorRT), STM32, quantization, optimisation inférence.}
    \item \small{\textbf{FPGA :} \textbf{VHDL}, Quartus Prime, Platform Designer (SoPC), ModelSim, Nios V.}
    \item \small{\textbf{Langages :} \textbf{Python} (NumPy, Pandas), \textbf{C/C++}, Bash.}
    \item \small{\textbf{DevOps / MLOps :} \textbf{Docker}, Git, CI/CD, Linux.}
\end{itemize}

% --- PROJETS ---
\section{Projets IA Embarquée \& Déploiement}
\begin{itemize}[leftmargin=*, label={}]

\item \textbf{Upgrade Jetson -- Déploiement IA sur GPU Embarqué} \hfill \textit{Challenge CoHoMa}
\vspace{-2pt}
\item \small{Défis de l'armée de terre Française : Déploiement et optimisation de modèles de vision sur Nvidia Jetson pour robotique autonome.}
    \begin{itemize}[leftmargin=0.4cm, nosep, label=\tiny$\bullet$]
        \item \small{\textbf{Déploiement GPU :} YOLOv10 sur Jetson, profilage inférence, optimisation mémoire.}
        \item \small{\textbf{Temps réel :} Contraintes latence pour détection d'objets embarquée.}
        \item \small{\textbf{Stack :} Python, C++, Docker, Linux.}
    \end{itemize}
    \vspace{3pt}

\item \textbf{Harmony Gloves -- TinyML sur Microcontrôleur} \hfill \textit{Laboratoire IOTA}
\vspace{-2pt}
\item \small{Modèle ML embarqué pour reconnaissance gestuelle temps réel.}
    \begin{itemize}[leftmargin=0.4cm, nosep, label=\tiny$\bullet$]
        \item \small{\textbf{Optimisation modèle :} Quantization pour déploiement sur STM32 (ressources limitées).}
        \item \small{\textbf{Acquisition :} Traitement signal capteurs IMU, fusion de données.}
    \end{itemize}
    \vspace{3pt}

\item \textbf{Système SoPC sur Cyclone V -- FPGA} \hfill \textit{TP ENSEA (En cours)}
\vspace{-2pt}
\item \small{Conception d'un système sur puce programmable avec soft-processeur.}
    \begin{itemize}[leftmargin=0.4cm, nosep, label=\tiny$\bullet$]
        \item \small{\textbf{Architecture :} Nios V via Platform Designer, configuration mémoire, bus Avalon.}
        \item \small{\textbf{Synthèse :} Flow Quartus complet, analyse timing, simulation ModelSim.}
    \end{itemize}
    \vspace{3pt}

\item \textbf{Classification Cancer du Poumon -- IndabaX 2025 (2\textsuperscript{e} Prix)} \hfill \textit{Hackathon IA}
\vspace{-2pt}
\item \small{Pipeline ML complet pour classification d'images médicales.}
    \begin{itemize}[leftmargin=0.4cm, nosep, label=\tiny$\bullet$]
        \item \small{\textbf{Modèles :} Random Forest (98.6\%), SVM -- comparaison et benchmarking.}
        \item \small{\textbf{Analyse :} Prétraitement images, évaluation performances, documentation.}
    \end{itemize}
    \vspace{3pt}

\item \textbf{Projet Taxi -- Simulation Temps Réel en C} \hfill \textit{Projet Académique}
\vspace{-2pt}
\item \small{Noyau de simulation multiprocessus sous Linux.}
    \begin{itemize}[leftmargin=0.4cm, nosep, label=\tiny$\bullet$]
        \item \small{\textbf{C/Linux :} Mémoire partagée, IPC, signaux UNIX, ordonnanceur.}
    \end{itemize}

\end{itemize}

% --- EXPÉRIENCE PROFESSIONNELLE ---
\section{Expériences Professionnelles}
\begin{itemize}[leftmargin=*, label={}]

    \resumeSubheading
      {Assistant de Recherche -- Deep Learning}{Juin 2025 -- Sept. 2025}
      {MILA (Institut Québécois d'IA) - Remote}{\textbf{Québec, Canada}}
      \begin{itemize}[leftmargin=0.4cm, nosep, label=\tiny$\bullet$]
        \item \small{Pipelines \textbf{PyTorch} pour entraînement et analyse de modèles (papier Grokking, OpenAI).}
        \item \small{Rigueur scientifique : reproduction expériences, documentation technique.}
      \end{itemize}
    \vspace{2pt}

    \resumeSubheading
      {Ingénieur Backend (Freelance)}{Oct. 2024 -- Août 2025}
      {CSC Fintech -- Architecture Bancaire}{Douala, Cameroun}
      \begin{itemize}[leftmargin=0.4cm, nosep, label=\tiny$\bullet$]
        \item \small{Architecture robuste (Docker, Redis, CI/CD), tests unitaires, revues de code.}
      \end{itemize}
\end{itemize}

% --- ENGAGEMENT ---
\section{Engagement \& Certifications}
\begin{itemize}[leftmargin=*, nosep]
    \item \textbf{Leadership :} Président Club Informatique, Chef Cellule Projet IA (Polytechnique Yaoundé).
    \item \textbf{Hackathons :} 2\textsuperscript{e} IndabaX Africa 2025, 1\textsuperscript{er} Mountain Hub, 1\textsuperscript{er} CITS National.
    \item \textbf{Certifications :} Supervised ML (DeepLearning.AI / Andrew Ng), Python for Data Science (IBM).
    \item \textbf{Langues :} Français (Natif), Anglais (Intermédiaire - Documentation technique).
\end{itemize}

\end{document}
